% Editable transcription of A4P1.pdf. Draft placeholders are intentional.
\documentclass[11pt,letterpaper]{article}
\usepackage{../styles/mae2250-handout}

\renewcommand{\HandoutNumber}{4}
\renewcommand{\HandoutTitle}{Coffee Grinder}
\renewcommand{\PartTitle}{Part 1: Sketching}
\renewcommand{\DueDate}{TBD}
\renewcommand{\TeamType}{Individual}

\begin{document}
\MakeHandoutHeader

For your last assignment, you will be prototyping a coffee grinder. You will manufacture a
test platform to test the rotating mechanisms, then finish with a CAD model of the final design.
Rather than using a power source, the coffee beans will be ground by hand. You will have the
following constraints:
\begin{itemize}
  \item There must be a place to both insert and retrieve the coffee beans (they may fall directly into
    a mug).
  \item You must use a set of gears available to you in the Taylor Design Studio (TDS).
  \item You must use the following set of burr gears to grind the coffee: \{link or image of burrs\}
  \item You may use a maximum of three parts machined in the Rapid Prototyping Lab (RPL) or
    Manufacturing Learning Studio (MLS) with Red Apron. Otherwise, all parts must be from
    McMaster or available in the TDS.
  \item You may not use adhesives or binding agents.
  \item You will have a budget of \$100 \{based on Sarah's and my BOMs\}.
  \item The coffee grinder must withstand pulling forces of 107N, and the beans must be ground
    with a torque of at least 1.5Nm.
  \item Your bounding box is 300mm*200mm*300mm.
\end{itemize}

\PartHeading{PART 1}
\begin{AssignmentSteps}
  \item \underline{Make sketches.}
    \begin{itemize}
      \item Make one sketch of only the internal mechanisms (shafts, gears, etc) and one sketch of the
        mechanisms plus the housing.
      \item On your housing sketch, indicate how the coffee beans enter and leave the coffee grinder.
    \end{itemize}
  \item \underline{Choose your bearings.}
    \begin{itemize}
      \item Make a free body diagram of the internal mechanisms of your coffee grinder.
        \begin{itemize}
          \item If it helps, also use a force flow or v-M diagram.
        \end{itemize}
      \item On your diagram, mark where you will need to place bearings. Write underneath what kind
        of bearings you will need (eg axial, thrust) and what loads they will need to support.
    \end{itemize}
\end{AssignmentSteps}

\textbf{Submit a document with your sketches, free body diagrams, and explanation of bearings to
Canvas by \{part 1 due date\}.}

Now that you have your initial sketches, you will go on to design and model your test structure.

\newpage

\textbf{\{Current intent of checklist setup is to make sure students are following good practices and
delivering on assignments while not punishing overlooked aspects of designs.\}}

\PartHeading{TA CHECKLIST}
\begin{Checklist}
  \item There is a sketch of the internal mechanisms
  \item There is a sketch of the full coffee grinder
  \item Sketches demonstrate understanding of how each part interacts with the others
  \item Sketches are clearly depicted and labelled
  \item FBD shows interactions from the hand, gears, and any supports listed in the sketch
  \item Coffee grinder will produce a given torque of 1.5 Nm at the listed force
  \item FBD (and any supporting diagrams) are clearly labelled
  \item Selected bearings can withstand proper forces
  \item Selected bearings are in the locations that see the highest forces
\end{Checklist}

\PartHeading{FEEDBACK CHECKLIST}
Launch point in case TAs aren't sure what to give feedback on:
\begin{Checklist}
  \item Feasibility -- Is this design possible in real life? Are there floating parts? Does the design
    look like it will tip over? Is the design overly complex?
  \item Manufacturability -- Will the student be able to assemble their coffee grinder with the tools
    provided? Are there parts that are impossible to machine? Are there parts that don't have a
    clear way of connecting to one another?
  \item Utility -- Can a person use this design? Is there a clear entrance and exit for the coffee beans?
    Is the housing stable when someone is using the crank?
\end{Checklist}

\end{document}
